\section{习题}

\subsection{本章速览}

\begin{itemize}
    \item 偏导数 \(f_x(x_0,y_0)\) 是固定 \(y=y_0\) 后对 \(x\) 求导，\(f_y(x_0,y_0)\) 是固定 \(x=x_0\) 后对 \(y\) 求导。它只刻画沿两个坐标方向的变化，偏导数存在既不能推出连续，也不能推出可微。
    \item 可微的定义是
    \[
        \varDelta z=f_x(x_0,y_0)\varDelta x+f_y(x_0,y_0)\varDelta y+o(\rho),
        \qquad \rho=\sqrt{\varDelta x^2+\varDelta y^2}.
    \]
    可微必连续、必可偏导；偏导数连续必可微。反向蕴含一般都不成立。
    \item 全微分为 \(\mathrm dz=f_x\,\mathrm dx+f_y\,\mathrm dy\)。当 \(\varDelta x,\varDelta y\) 很小时，
    \[
        f(x_0+\varDelta x,y_0+\varDelta y)\approx f(x_0,y_0)+f_x(x_0,y_0)\varDelta x+f_y(x_0,y_0)\varDelta y.
    \]
    \item 复合函数求导先理清变量依赖关系，再逐层使用链式法则；二阶求导时，一阶偏导仍是中间变量的函数，必须再次展开。
    \item 方程 \(F(x,y,z)=0\) 在 \(F_z\neq0\) 时确定隐函数 \(z=z(x,y)\)，且
    \[
        \frac{\partial z}{\partial x}=-\frac{F_x}{F_z},\qquad
        \frac{\partial z}{\partial y}=-\frac{F_y}{F_z}.
    \]
    \item 方向导数 \(\dfrac{\partial f}{\partial \mathbf l}=\nabla f\cdot\mathbf e\)，其中 \(\mathbf e\) 是单位方向向量；当 \(\nabla f\neq\mathbf 0\) 时，沿梯度方向方向导数最大，最大值为 \(|\nabla f|\)。极值问题先用驻点与 Hessian 判别，条件极值用 Lagrange 乘数法。
\end{itemize}

\subsection{基础过关题}

\textbf{1.} \textbf{（概念辨析）} 判断下列说法是否正确，并说明理由：
(1) 若 \(f(x,y)\) 在点 \(P_0\) 可微，则 \(f\) 在 \(P_0\) 连续，且 \(f_x(P_0),f_y(P_0)\) 都存在；
(2) 若 \(f_x(P_0),f_y(P_0)\) 都存在，则 \(f\) 在 \(P_0\) 连续；
(3) 若 \(f\) 在 \(P_0\) 可微，则对任一单位向量 \(\mathbf e\)，方向导数 \(\dfrac{\partial f}{\partial\mathbf e}(P_0)\) 存在，且等于 \(\nabla f(P_0)\cdot\mathbf e\)；
(4) 若 \(f_{xy}(P_0)\) 与 \(f_{yx}(P_0)\) 都存在，则二者相等。

\ifshowsolutions%
\boxed{\textbf{解}}
\textbf{思路：} 逐条对照“可微、连续、可偏导、偏导连续”之间的蕴含关系，假命题举反例。

(1) 正确。可微的定义本身给出
\[
\varDelta f=A\varDelta x+B\varDelta y+o(\rho),
\]
当 \(\rho\to0\) 时 \(\varDelta f\to0\)，故连续；又由可微的必要条件，\(A=f_x(P_0)\)，\(B=f_y(P_0)\)，所以两个偏导数存在。

(2) 错误。取
\[
f(x,y)=\begin{cases}
\dfrac{xy}{x^2+y^2},&(x,y)\neq(0,0),\\[6pt]
0,&(x,y)=(0,0).
\end{cases}
\]
由定义
\[
f_x(0,0)=\lim_{h\to0}\frac{f(h,0)-f(0,0)}{h}=0,
\]
同理 \(f_y(0,0)=0\)。但沿直线 \(y=x\) 有 \(f(x,x)=\dfrac12\to\dfrac12\neq f(0,0)\)，故 \(f\) 在原点不连续。

(3) 正确。设 \(\mathbf e=(\cos\alpha,\sin\alpha)\)，令 \(\varphi(t)=f(P_0+t\mathbf e)\)。由可微性，
\[
\varphi'(0)=f_x(P_0)\cos\alpha+f_y(P_0)\sin\alpha=\nabla f(P_0)\cdot\mathbf e.
\]
这正是方向导数。

(4) 错误。取
\[
f(x,y)=\begin{cases}
\dfrac{xy(x^2-y^2)}{x^2+y^2},&(x,y)\neq(0,0),\\[6pt]
0,&(x,y)=(0,0).
\end{cases}
\]
当 \(y\neq0\) 时，
\[
f_x(0,y)=\lim_{h\to0}\frac{f(h,y)-f(0,y)}{h}
=\lim_{h\to0}\frac{y(h^2-y^2)}{h^2+y^2}=-y,
\]
且 \(f_x(0,0)=0\)，故
\[
f_{xy}(0,0)=\lim_{y\to0}\frac{f_x(0,y)-f_x(0,0)}{y}=-1.
\]
同理，当 \(x\neq0\) 时，
\[
f_y(x,0)=\lim_{k\to0}\frac{f(x,k)-f(x,0)}{k}
=\lim_{k\to0}\frac{x(x^2-k^2)}{x^2+k^2}=x,
\]
故
\[
f_{yx}(0,0)=\lim_{x\to0}\frac{f_y(x,0)-f_y(0,0)}{x}=1.
\]
两者不相等。保证 \(f_{xy}=f_{yx}\) 的常用条件是两个混合偏导数在该点连续，仅有“存在”不够。

\boxed{\text{答案：(1)(3) 正确；(2)(4) 错误。}}
\else \answerarea[8] \fi%

\textbf{2.} \textbf{（全微分与近似计算）} 设 \(z=\sqrt{x^2+y^2}\)。

(1) 求 \(\mathrm dz\big|_{(3,4)}\)；

(2) 用全微分近似计算 \(\sqrt{2.98^2+4.03^2}\)，并说明近似的阶。

\ifshowsolutions%
\boxed{\textbf{解}}
\textbf{思路：} 先求两个偏导数，再套全微分公式；近似时把待求点写成 \((x_0+\varDelta x,y_0+\varDelta y)\)。

(1)
\[
\frac{\partial z}{\partial x}=\frac{x}{\sqrt{x^2+y^2}},\qquad
\frac{\partial z}{\partial y}=\frac{y}{\sqrt{x^2+y^2}}.
\]
在 \((3,4)\) 处，\(\sqrt{x^2+y^2}=5\)，故
\[
z_x(3,4)=\frac35,\qquad z_y(3,4)=\frac45,
\]
从而
\[
\mathrm dz\big|_{(3,4)}=\frac35\,\mathrm dx+\frac45\,\mathrm dy.
\]

(2) 取 \(x_0=3,\ y_0=4\)，则 \(\varDelta x=-0.02,\ \varDelta y=0.03\)，且 \(z(3,4)=5\)。于是
\[
\varDelta z\approx \mathrm dz\big|_{(3,4)}
=\frac35(-0.02)+\frac45(0.03)
=-0.012+0.024=0.012.
\]
所以
\[
\sqrt{2.98^2+4.03^2}\approx 5+0.012=5.012.
\]
真值为 \(5.012115\ldots\)，一阶近似的误差约为 \(1.2\times10^{-4}\)，是 \(|\varDelta x|,|\varDelta y|\) 的二阶小量。

\boxed{\text{答案：(1) $\mathrm dz\big|_{(3,4)}=\dfrac35\mathrm dx+\dfrac45\mathrm dy$；(2) 约为 $5.012$。}}
\else \answerarea[8] \fi%

\textbf{3.} \textbf{（复合函数的一阶偏导）} 设 \(f(u,v)\) 具有连续偏导数，\(z=f\left(xy,\dfrac{x}{y}\right)\)（\(y\neq0\)）。
(1) 求 \(\dfrac{\partial z}{\partial x},\dfrac{\partial z}{\partial y}\)；
(2) 若 \(f_u\left(2,\dfrac12\right)=3,\ f_v\left(2,\dfrac12\right)=-1\)，求 \(z_x(1,2),z_y(1,2)\)。

\ifshowsolutions%
\boxed{\textbf{解}}
\textbf{思路：} 令 \(u=xy,\ v=\dfrac{x}{y}\)，先对中间变量用链式法则，再代入具体数值。

(1) 因为
\[
u_x=y,\quad u_y=x,\quad v_x=\frac1y,\quad v_y=-\frac{x}{y^2},
\]
所以
\[
\frac{\partial z}{\partial x}=y f_u+\frac1y f_v,
\qquad
\frac{\partial z}{\partial y}=x f_u-\frac{x}{y^2}f_v.
\]

(2) 当 \((x,y)=(1,2)\) 时，\(u=2,\ v=\dfrac12\)。代入上式：
\[
z_x(1,2)=2\cdot3+\frac12\cdot(-1)=\frac{11}{2},
\]
\[
z_y(1,2)=1\cdot3-\frac{1}{4}\cdot(-1)=\frac{13}{4}.
\]

\boxed{\text{答案：(1) $z_x=yf_u+\dfrac1y f_v,\ z_y=xf_u-\dfrac{x}{y^2}f_v$；(2) $z_x(1,2)=\dfrac{11}{2},\ z_y(1,2)=\dfrac{13}{4}$。}}
\else \answerarea[8] \fi%

\textbf{4.} \textbf{（隐函数求导）} 设 \(z=z(x,y)\) 由方程
\[
x^2+y^2+z^2=6
\]
在点 \((1,1,2)\) 附近确定。求 \(\dfrac{\partial z}{\partial x},\dfrac{\partial z}{\partial y}\) 在 \((1,1)\) 处的值，并求 \(\dfrac{\partial^2 z}{\partial x^2},\dfrac{\partial^2 z}{\partial x\partial y}\) 在 \((1,1)\) 处的值。

\ifshowsolutions%
\boxed{\textbf{解}}
\textbf{思路：} 先验证隐函数存在定理的条件，再对恒等式两边求导；求二阶导时，把 \(z_x,z_y\) 仍看作 \(x,y\) 的函数继续求导。

令 \(F(x,y,z)=x^2+y^2+z^2-6\)。在 \((1,1,2)\) 处
\[
F_z=2z=4\neq0,
\]
故方程在该点附近确定可微隐函数 \(z=z(x,y)\)。

对 \(F(x,y,z(x,y))=0\) 两边分别关于 \(x,y\) 求偏导：
\[
2x+2zz_x=0,\qquad 2y+2zz_y=0,
\]
于是
\[
z_x=-\frac{x}{z},\qquad z_y=-\frac{y}{z}.
\]
在 \((1,1,2)\) 处，
\[
z_x(1,1)=-\frac12,\qquad z_y(1,1)=-\frac12.
\]

再对 \(2x+2zz_x=0\) 关于 \(x\) 求偏导：
\[
2+2z_x^2+2zz_{xx}=0,
\]
故
\[
z_{xx}=-\frac{1+z_x^2}{z}.
\]
代入 \(z=2,\ z_x=-\dfrac12\)，得
\[
z_{xx}(1,1)=-\frac{1+\frac14}{2}=-\frac58.
\]

对 \(2x+2zz_x=0\) 关于 \(y\) 求偏导：
\[
2z_yz_x+2zz_{xy}=0,
\]
故
\[
z_{xy}=-\frac{z_xz_y}{z}.
\]
代入 \(z_x=z_y=-\dfrac12,\ z=2\)，得
\[
z_{xy}(1,1)=-\frac{\frac14}{2}=-\frac18.
\]

\boxed{\text{答案：$z_x(1,1)=z_y(1,1)=-\dfrac12,\ z_{xx}(1,1)=-\dfrac58,\ z_{xy}(1,1)=-\dfrac18$。}}
\else \answerarea[10] \fi%

\textbf{5.} \textbf{（方向导数与梯度）} 设 \(f(x,y)=x^2y+y^3\)，点 \(P_0(1,2)\)。
(1) 求 \(f\) 在 \(P_0\) 处沿从 \(P_0\) 指向 \((3,3)\) 的方向的方向导数；
(2) 求 \(f\) 在 \(P_0\) 处的最大方向导数及取得最大方向导数的方向；
(3) 求 \(f\) 在 \(P_0\) 处方向导数为零的方向。

\ifshowsolutions%
\boxed{\textbf{解}}
\textbf{思路：} 先求梯度，再把方向化为单位向量，利用 \(\dfrac{\partial f}{\partial\mathbf l}=\nabla f\cdot\mathbf e\)。

先求偏导：
\[
f_x=2xy,\qquad f_y=x^2+3y^2.
\]
在 \(P_0(1,2)\) 处，
\[
\nabla f(1,2)=(4,13).
\]

(1) 方向向量为 \((3,3)-(1,2)=(2,1)\)，单位化为
\[
\mathbf e=\frac{1}{\sqrt5}(2,1).
\]
故
\[
\frac{\partial f}{\partial\mathbf l}\bigg|_{P_0}
=(4,13)\cdot\frac{(2,1)}{\sqrt5}
=\frac{8+13}{\sqrt5}
=\frac{21}{\sqrt5}
=\frac{21\sqrt5}{5}.
\]

(2) 方向导数的最大值为
\[
|\nabla f(1,2)|=\sqrt{4^2+13^2}=\sqrt{185},
\]
取得最大方向导数的方向为梯度方向
\[
\frac{\nabla f(1,2)}{|\nabla f(1,2)|}
=\frac{1}{\sqrt{185}}(4,13).
\]

(3) 方向导数为零当且仅当方向与梯度垂直。取
\[
\mathbf e=\frac{1}{\sqrt{185}}(-13,4)
\quad\text{或}\quad
\mathbf e=\frac{1}{\sqrt{185}}(13,-4),
\]
此时 \((4,13)\cdot(-13,4)=-52+52=0\)，方向导数为零。

\boxed{\text{答案：(1) $\dfrac{21\sqrt5}{5}$；(2) 最大值 $\sqrt{185}$，方向 $\dfrac{(4,13)}{\sqrt{185}}$；(3) 方向 $\pm\dfrac{(-13,4)}{\sqrt{185}}$。}}
\else \answerarea[8] \fi%

\subsection{综合提高题}

\textbf{6.} \textbf{（偏导、连续、方向导数与可微性）} 设
\[
f(x,y)=\begin{cases}
\dfrac{x^2y}{x^2+y^2},&(x,y)\neq(0,0),\\[8pt]
0,&(x,y)=(0,0).
\end{cases}
\]
(1) 求 \(f_x(0,0),f_y(0,0)\)；
(2) 讨论 \(f\) 在原点处的连续性；
(3) 求 \(f\) 在原点处沿任意单位方向 \(\mathbf e=(\cos\theta,\sin\theta)\) 的方向导数；
(4) 判断 \(f\) 在原点处是否可微，并说明“各方向方向导数都存在”与“可微”的区别。

\ifshowsolutions%
\boxed{\textbf{解}}
\textbf{思路：} 分段函数在分界点处一律回到定义；方向导数用射线 \(t\mapsto(t\cos\theta,t\sin\theta)\) 计算；可微性用
\[
\frac{\varDelta f-f_x(0,0)\varDelta x-f_y(0,0)\varDelta y}{\rho}\to0
\]
是否成立来判断。

(1) 因为 \(f(h,0)=0,\ f(0,k)=0\)，由定义
\[
f_x(0,0)=\lim_{h\to0}\frac{f(h,0)-f(0,0)}{h}=0,
\]
\[
f_y(0,0)=\lim_{k\to0}\frac{f(0,k)-f(0,0)}{k}=0.
\]

(2) 当 \((x,y)\neq(0,0)\) 时，
\[
|f(x,y)|=\frac{x^2|y|}{x^2+y^2}\le |y|,
\]
因为 \(\dfrac{x^2}{x^2+y^2}\le1\)。于是 \((x,y)\to(0,0)\) 时 \(f(x,y)\to0=f(0,0)\)，故 \(f\) 在原点连续。

(3) 沿方向 \(\mathbf e=(\cos\theta,\sin\theta)\)，
\[
f(t\cos\theta,t\sin\theta)
=\frac{t^2\cos^2\theta\cdot t\sin\theta}{t^2}
=t\cos^2\theta\sin\theta\qquad(t\neq0).
\]
按方向导数定义，
\[
\frac{\partial f}{\partial\mathbf e}(0,0)
=\lim_{t\to0}\frac{f(t\cos\theta,t\sin\theta)-f(0,0)}{t}
=\cos^2\theta\sin\theta.
\]
因此沿任意方向的方向导数都存在（例如沿 \(\theta=0\) 为 \(0\)，沿 \(\theta=\frac{\pi}{4}\) 为 \(\frac{\sqrt2}{4}\)）。

(4) 若 \(f\) 在原点可微，由于 \(f_x(0,0)=f_y(0,0)=0\)，应有
\[
\frac{\varDelta f}{\rho}\to0.
\]
但取 \(\varDelta x=\rho\cos\theta,\ \varDelta y=\rho\sin\theta\)（\(\rho>0\)），则
\[
\frac{\varDelta f}{\rho}
=\frac{\rho^3\cos^2\theta\sin\theta/\rho^2}{\rho}
=\cos^2\theta\sin\theta,
\]
其值随方向 \(\theta\) 改变（例如 \(\theta=0\) 时为 \(0\)，\(\theta=\frac{\pi}{4}\) 时为 \(\frac{\sqrt2}{4}\)），并不趋于唯一的 \(0\)。故 \(f\) 在原点不可微。

这说明：即使沿所有方向的方向导数都存在，只要这些方向导数不能统一写成“梯度与单位方向的内积”这一线性形式，函数仍不可微。本题中方向导数为 \(\cos^2\theta\sin\theta\)，它不是 \((\cos\theta,\sin\theta)\) 的线性函数。

\boxed{\text{答案：(1) $f_x(0,0)=f_y(0,0)=0$；(2) 连续；(3) $\cos^2\theta\sin\theta$；(4) 不可微。}}
\else \answerarea[10] \fi%

\textbf{7.} \textbf{（二阶复合求导与调和函数）} 设 \(f(u,v)\) 具有二阶连续偏导数，\(z=f(x^2-y^2,\,2xy)\)。
(1) 求 \(\dfrac{\partial z}{\partial x},\dfrac{\partial z}{\partial y}\)；
(2) 求 \(\dfrac{\partial^2 z}{\partial x^2}+\dfrac{\partial^2 z}{\partial y^2}\)，并用 \(f\) 的偏导数表示；
(3) 若 \(f\) 满足 \(f_{uu}+f_{vv}=0\)，判断 \(z\) 是否也满足 \(z_{xx}+z_{yy}=0\)。

\ifshowsolutions%
\boxed{\textbf{解}}
\textbf{思路：} 令 \(u=x^2-y^2,\ v=2xy\)。一阶求导后，\(f_u,f_v\) 仍是 \((u,v)\) 的函数，二阶求导必须再次使用链式法则。以下 \(f_u,f_v,f_{uu},f_{uv},f_{vv}\) 均在 \((u,v)\) 处取值，且利用 \(f_{uv}=f_{vu}\)。

(1)
\[
z_x=f_u\cdot2x+f_v\cdot2y=2xf_u+2yf_v,
\]
\[
z_y=f_u\cdot(-2y)+f_v\cdot2x=-2yf_u+2xf_v.
\]

(2) 对 \(z_x\) 再关于 \(x\) 求偏导：
\[
\begin{aligned}
z_{xx}
&=2f_u+2x\bigl(2xf_{uu}+2yf_{uv}\bigr)
 +2y\bigl(2xf_{vu}+2yf_{vv}\bigr)\\
&=2f_u+4x^2f_{uu}+8xyf_{uv}+4y^2f_{vv}.
\end{aligned}
\]
对 \(z_y\) 再关于 \(y\) 求偏导：
\[
\begin{aligned}
z_{yy}
&=-2f_u-2y\bigl(-2yf_{uu}+2xf_{uv}\bigr)
 +2x\bigl(-2yf_{vu}+2xf_{vv}\bigr)\\
&=-2f_u+4y^2f_{uu}-8xyf_{uv}+4x^2f_{vv}.
\end{aligned}
\]
两式相加，\(f_u\) 项与 \(f_{uv}\) 项相消，得
\[
z_{xx}+z_{yy}
=4(x^2+y^2)\bigl(f_{uu}+f_{vv}\bigr).
\]

(3) 若 \(f_{uu}+f_{vv}=0\)，则由 (2) 立即得到
\[
z_{xx}+z_{yy}=0.
\]
即调和函数在变换 \((u,v)=(x^2-y^2,2xy)\) 下仍为调和函数。该变换正是复映射 \(w=(x+iy)^2\) 的实部与虚部。

\boxed{\text{\parbox{0.9\linewidth}{\raggedright 答案：(1) $z_x=2xf_u+2yf_v,\ z_y=-2yf_u+2xf_v$；(2) $z_{xx}+z_{yy}=4(x^2+y^2)(f_{uu}+f_{vv})$；(3) 满足。}}}
\else \answerarea[10] \fi%

\textbf{8.} \textbf{（无条件极值与 Hessian 判别）} 求函数
\[
f(x,y)=x^4+y^4-4xy
\]
的驻点，写出 Hessian 矩阵，判定每个驻点的类型，并求出极值。

\ifshowsolutions%
\boxed{\textbf{解}}
\textbf{思路：} 先解方程组 \(f_x=f_y=0\) 得驻点，再在每个驻点计算 Hessian 矩阵，用 \(AC-B^2\) 与 \(A\) 的符号判别。

求一阶偏导：
\[
f_x=4x^3-4y,\qquad f_y=4y^3-4x.
\]
令其为零：
\[
y=x^3,\qquad x=y^3.
\]
代入得 \(x=x^9\)，即
\[
x(x^8-1)=0.
\]
所以 \(x=0,1,-1\)，对应 \(y=0,1,-1\)。三个驻点为
\[
P_0(0,0),\qquad P_1(1,1),\qquad P_2(-1,-1).
\]

再求二阶偏导：
\[
f_{xx}=12x^2,\qquad f_{xy}=-4,\qquad f_{yy}=12y^2.
\]
Hessian 矩阵为
\[
H(x,y)=
\begin{pmatrix}
12x^2 & -4\\
-4 & 12y^2
\end{pmatrix}.
\]

在 \(P_0(0,0)\) 处，
\[
H(0,0)=\begin{pmatrix}0&-4\\-4&0\end{pmatrix},
\qquad
AC-B^2=0\cdot0-(-4)^2=-16<0,
\]
Hessian 不定，故 \(P_0\) 是鞍点，不是极值点。

在 \(P_1(1,1)\) 与 \(P_2(-1,-1)\) 处，
\[
H=\begin{pmatrix}12&-4\\-4&12\end{pmatrix},
\qquad
AC-B^2=144-16=128>0,
\]
且 \(A=12>0\)，Hessian 正定，故两点都是严格极小值点。极小值为
\[
f(1,1)=1+1-4=-2,\qquad f(-1,-1)=1+1-4=-2.
\]

又当 \(x^2+y^2\to\infty\) 时 \(f(x,y)\to+\infty\)，所以这两个极小值也是全局最小值，函数没有最大值。

\boxed{\text{\parbox{0.9\linewidth}{\raggedright 答案：驻点为 $(0,0),(1,1),(-1,-1)$；$(0,0)$ 为鞍点；$(1,1),(-1,-1)$ 为极小值点，极小值均为 $-2$，且为全局最小值。}}}
\else \answerarea[10] \fi%

\textbf{9.} \textbf{（条件极值·Lagrange 乘数法）} 求函数
\[
f(x,y,z)=xy+yz+zx
\]
在单位球面 \(x^2+y^2+z^2=1\) 上的最大值与最小值，并指出最大值点。

\ifshowsolutions%
\boxed{\textbf{解}}
\textbf{思路：} 单位球面是有界闭集，\(f\) 连续，最值一定存在。构造 Lagrange 函数，解方程组后比较各候选值；注意按 \(x+y+z\) 是否为零分类。

令
\[
L=xy+yz+zx-\lambda(x^2+y^2+z^2-1).
\]
拉格朗日方程组为
\[
\begin{cases}
y+z=2\lambda x,\\
x+z=2\lambda y,\\
x+y=2\lambda z,\\
x^2+y^2+z^2=1.
\end{cases}
\]
将前三式相加：
\[
2(x+y+z)=2\lambda(x+y+z),
\]
即
\[
(\lambda-1)(x+y+z)=0.
\]

\textbf{情形一：}\(\lambda=1\)。此时前两式相减：
\[
(y+z)-(x+z)=2(x-y),
\]
即 \(y-x=2x-2y\)，所以 \(3(y-x)=0\)，得 \(x=y\)。同理 \(y=z\)。于是
\[
x=y=z=\pm\frac1{\sqrt3},
\]
此时
\[
f=3\cdot\frac13=1.
\]

\textbf{情形二：}\(x+y+z=0\)。将三个方程分别乘以 \(x,y,z\) 后相加：
\[
2(xy+yz+zx)=2\lambda(x^2+y^2+z^2)=2\lambda,
\]
故 \(\lambda=f\)。又由 \(x+y+z=0\) 与 \(x^2+y^2+z^2=1\)，
\[
0=(x+y+z)^2=1+2(xy+yz+zx)=1+2f,
\]
所以
\[
f=-\frac12.
\]
例如 \((x,y,z)=\left(\frac1{\sqrt2},-\frac1{\sqrt2},0\right)\) 满足条件，确实取得该值。

比较两种情形，最大值为 \(1\)，最小值为 \(-\dfrac12\)。

\boxed{\text{\parbox{0.9\linewidth}{\raggedright 答案：最大值为 $1$，在 $\pm\left(\dfrac1{\sqrt3},\dfrac1{\sqrt3},\dfrac1{\sqrt3}\right)$ 处取得；最小值为 $-\dfrac12$，在圆 $x+y+z=0,\ x^2+y^2+z^2=1$ 上取得。}}}
\else \answerarea[10] \fi%

\subsection{真题风格与综合训练}

\textbf{10.} \textbf{（真题风格·选择题）} 设 \(f(x,y)\) 在点 \(P_0\) 处可微，且
\[
\nabla f(P_0)=(2,-1).
\]
则下列结论中正确的是（\quad）

\begin{enumerate}[label=(\Alph*),itemsep=2pt,topsep=2pt]
    \item \(f\) 在 \(P_0\) 处沿方向 \((1,2)\) 的方向导数为 \(0\)
    \item \(f\) 在 \(P_0\) 处沿方向 \((-2,1)\) 的方向导数为 \(\sqrt5\)
    \item \(f\) 在 \(P_0\) 处的最大方向导数为 \(1\)
    \item \(f\) 在 \(P_0\) 处沿任意方向的方向导数都存在，但未必等于 \(\nabla f(P_0)\) 与该方向单位向量的内积
\end{enumerate}

\ifshowsolutions%
\boxed{\textbf{解}}
\textbf{思路：} 可微时，沿单位向量 \(\mathbf e\) 的方向导数等于梯度与 \(\mathbf e\) 的内积；最大方向导数为梯度的模。

(A) 单位方向为 \(\dfrac{(1,2)}{\sqrt5}\)，方向导数为
\[
(2,-1)\cdot\frac{(1,2)}{\sqrt5}=\frac{2-2}{\sqrt5}=0,
\]
正确。

(B) 方向导数为
\[
(2,-1)\cdot\frac{(-2,1)}{\sqrt5}=\frac{-4-1}{\sqrt5}=-\sqrt5\neq\sqrt5,
\]
错误。

(C) 最大方向导数为
\[
|\nabla f(P_0)|=\sqrt{2^2+(-1)^2}=\sqrt5\neq1,
\]
错误。

(D) 可微性保证方向导数公式成立，即一定等于梯度与该方向单位向量的内积，所以“未必等于”错误。

\boxed{\textbf{答案：}(A)}
\else \answerarea[4] \fi%

\textbf{11.} \textbf{（真题风格·填空题）} 设
\[
f(x,y)=x^2e^y+y\ln x\qquad(x>0),
\]
则 \(f_{xy}(1,0)=\underline{\hspace{2cm}}\)，\(f\) 在点 \((1,0)\) 处的 Hessian 矩阵的行列式为 \(\underline{\hspace{2cm}}\)。

\ifshowsolutions%
\boxed{\textbf{解}}
\textbf{思路：} 先求一阶偏导，再求二阶偏导；Hessian 矩阵为
\(\begin{pmatrix}f_{xx}&f_{xy}\\ f_{xy}&f_{yy}\end{pmatrix}\)。

一阶偏导：
\[
f_x=2xe^y+\frac{y}{x},\qquad
f_y=x^2e^y+\ln x.
\]
二阶偏导：
\[
f_{xx}=2e^y-\frac{y}{x^2},\qquad
f_{xy}=2xe^y+\frac1x,\qquad
f_{yy}=x^2e^y.
\]
在 \((1,0)\) 处，
\[
f_{xx}=2,\qquad f_{xy}=2+1=3,\qquad f_{yy}=1.
\]
故 Hessian 矩阵为
\[
H(1,0)=\begin{pmatrix}2&3\\3&1\end{pmatrix},
\]
其行列式为
\[
\det H(1,0)=2\cdot1-3\cdot3=-7.
\]

\boxed{\text{答案：$3$；$-7$。}}
\else \answerarea[4] \fi%

\textbf{12.} \textbf{（真题风格·解答题）} 求函数
\[
f(x,y)=x^2+y^2+xy-x-y+1
\]
在闭区域
\[
D=\{(x,y)\mid x^2+y^2\le2\}
\]
上的最大值与最小值。

\ifshowsolutions%
\boxed{\textbf{解}}
\textbf{思路：} 有界闭区域上的连续函数必有最值。先求区域内部的驻点，再求边界 \(x^2+y^2=2\) 上的最值，最后比较。

\textbf{（1）内部驻点。}
\[
f_x=2x+y-1,\qquad f_y=2y+x-1.
\]
令其为零：
\[
\begin{cases}
2x+y=1,\\
x+2y=1.
\end{cases}
\]
两式相减得 \(x=y\)，再代入得 \(x=y=\dfrac13\)，且
\[
\left(\frac13\right)^2+\left(\frac13\right)^2=\frac29<2,
\]
所以该点位于 \(D\) 的内部。此时
\[
f\left(\frac13,\frac13\right)
=\frac19+\frac19+\frac19-\frac13-\frac13+1
=\frac23.
\]
又
\[
f_{xx}=2,\qquad f_{xy}=1,\qquad f_{yy}=2,
\]
Hessian 矩阵 \(\begin{pmatrix}2&1\\1&2\end{pmatrix}\) 正定，故该点是内部极小值点。

\textbf{（2）边界上的最值。} 在边界 \(x^2+y^2=2\) 上，令
\[
x=\sqrt2\cos\theta,\qquad y=\sqrt2\sin\theta.
\]
则
\[
f=2+xy-x-y+1=3+xy-(x+y).
\]
而
\[
xy=2\sin\theta\cos\theta=\sin2\theta,\qquad
x+y=\sqrt2(\cos\theta+\sin\theta).
\]
记 \(s=\cos\theta+\sin\theta\)，则 \(s^2=1+\sin2\theta\)，即 \(\sin2\theta=s^2-1\)。于是
\[
f=3+(s^2-1)-\sqrt2s
=s^2-\sqrt2s+2
=\left(s-\frac{\sqrt2}{2}\right)^2+\frac32.
\]
因为 \(s=\cos\theta+\sin\theta\in[-\sqrt2,\sqrt2]\)，所以
\[
\min_{\text{边界}}f=\frac32
\quad\left(s=\frac{\sqrt2}{2}\right),
\]
\[
\max_{\text{边界}}f
=\left(-\sqrt2-\frac{\sqrt2}{2}\right)^2+\frac32
=\frac92+\frac32=6
\quad(s=-\sqrt2).
\]
当 \(s=-\sqrt2\) 时，\(\cos\theta=\sin\theta=-\dfrac1{\sqrt2}\)，即 \(x=y=-1\)，故最大值在 \((-1,-1)\) 处取得。

\textbf{（3）比较。} 内部极小值 \(\dfrac23\) 小于边界最小值 \(\dfrac32\)，而边界最大值为 \(6\)。因此
\[
\min_D f=\frac23,\qquad \max_D f=6.
\]

\boxed{\text{答案：最小值为 $\dfrac23$，在 $\left(\dfrac13,\dfrac13\right)$ 处取得；最大值为 $6$，在 $(-1,-1)$ 处取得。}}
\else \answerarea[10] \fi%
